% Generated draft: requires mathematical and source review.
\documentclass[11pt]{article}
\usepackage[margin=25mm]{geometry}
\usepackage{fontspec}
\setmainfont{Latin Modern Roman}
\usepackage{amsmath,amssymb,mathtools,unicode-math}
\setmathfont{Latin Modern Math}
\usepackage{enumitem,needspace,xurl,hyperref}
\setlength{\parindent}{0pt}
\setlength{\parskip}{5pt}
\setlength{\emergencystretch}{3em}
\newcommand{\N}{\mathbb{N}}
\newcommand{\Z}{\mathbb{Z}}
\newcommand{\Q}{\mathbb{Q}}
\newcommand{\R}{\mathbb{R}}
\newcommand{\C}{\mathbb{C}}
\newcommand{\F}{\mathbb{F}}
\newcommand{\A}{\mathbb{A}}
\newcommand{\PP}{\mathbb{P}}
\newcommand{\E}{\mathbb{E}}
\newcommand{\eps}{\varepsilon}
\newcommand{\dd}{\,\mathrm{d}}
\begin{document}
\section*{UnsolvedMath ID: 30000467}
\textbf{Source problem code:} catalogue entry
\subsection{Definitions and mathematical statement}
An $\R$-tree is a metric space in which any two points are joined by a unique arc and this arc is isometric to a real interval. Use the usual convention that a group action on an $\R$-tree is by isometries. Such an action is minimal if there is no proper nonempty invariant subtree, and has no global fixed point if
\[
 \nexists x\in T\ \forall g\in G:\quad gx=x.
\]
The source's precise subtree convention, including any closedness requirement, must be checked.

Classify or otherwise describe the finitely generated groups $G$ for which there exist an $\R$-tree $T$ and an action $G\curvearrowright T$ that is both minimal and without a global fixed point. In symbols, the class of groups to be described is
\[
 \{G:G\text{ is finitely generated and admits such an action on some }T\}.
\]
A group-isomorphism classification or an intrinsic necessary-and-sufficient structural description would address the stated target; the source does not prescribe the form of the answer. It asks about the groups admitting actions, not explicitly for classification of all actions up to equivariant isometry.

No finite-presentation hypothesis is supplied. Nor does the record require the action to be faithful, free, stable, small, or to have abelian, trivial, or otherwise controlled arc stabilizers. The tree need not be simplicial, locally compact, or complete by any explicit assumption in the extraction. These choices materially affect structure theorems, so none is silently imported. The supplied assessment mentions Rips-type special cases without giving their hypotheses; they cannot replace the unrestricted requested class.
\subsection{Short English statement}
The question asks which finitely generated groups admit a minimal isometric action on a real tree with no common fixed point. It seeks a structural description of the groups, with no finite-presentation or stabilizer restrictions supplied.
\subsection{Sources}
\begin{itemize}
\item Local UnsolvedMath registry, permanent ID 30000467, code catalogue entry, ``Groups Acting Minimally on Real Trees.'' Proposed year 2006; proposer unspecified.
\item \emph{Topological and Geometric Methods in Group Theory} (2007), Oberwolfach Reports, pp.\ 1169--1214, DOI 10.4171/owr/2006/19; \url{https://doi.org/10.4171/owr/2006/19}. The displayed report year 2007 differs from the DOI and proposed year 2006.
\item Statement status corrected\_verified; the supplied verification says official-report content was checked and notation, context, and question form normalized without changing the claim.
\item Literature triage checked 2026-08-21: open. It calls this a broad classification problem, reports substantial Rips-type special cases, and says they do not provide the unrestricted classification. No separate theorem citation or hypotheses are supplied. Created 2026-08-25; updated 2026-09-16. Current status and cited claims require independent verification.
\item Registry curation by Ulam AI / UnsolvedMath Contributors, CC BY 4.0; the underlying report retains its supplied provenance. Unreviewed statement-only draft.
\end{itemize}
\end{document}
